\section{OpenAI、misalignment報告フレームワークを公開}
OpenAIは、モデルのmisalignment事例を追跡・調査・公開する基準と期限を定めた新フレームワークを発表し、過去6か月に観測した6件の事例報告も同時公開した。未解明・未緩和の段階でも公開する方針を示し、今後の経験と公開フィードバックに応じて改訂するとしている。
\dailyxsource{OpenAI (@OpenAI)}{https://x.com/OpenAI/status/2100344867507327087}

\section{OpenAI、Astra for Lawを発表}
GPT-6 Astraを法務分析・執筆向けに構成した「Astra for Law」が発表された。米国判例・法令等を対象とするLegal Search Indexを組み合わせ、初期は選定法律事務所向けTrusted AccessとしてChatGPTとCodexで提供、APIは近日予定とされた。ベンチマーク値はOpenAI自己評価である。
\dailyxsource{OpenAI (@OpenAI)}{https://x.com/OpenAI/status/2100679992720142459}

\section{Anthropic、AI開発ペースを測る3指標を公開}
Anthropicは、自社R\&DにAIがどれだけ関与するか、エージェント監督の質、計算資源配分という3指標を公開した。2026年8月時点の社内スナップショットでは「AI leads」が26\%、AL3以上が90\%超、最頻内部プラットフォーム上で同時約3万エージェントと報告している。数値は同社自己測定である。
\dailyxsource{Anthropic (@AnthropicAI)}{https://x.com/AnthropicAI/status/2100684274114699295}

\section{生命科学向けVerification Programを開始}
Anthropicは生命科学専門家向けLife Sciences Verification Programをベータ公開した。より広い生物学作業を許可するセーフガード付き環境でMythos、Opus、Sonnetを利用でき、Mythosを含む初の一般応募枠と説明している。high-risk経路では提供モデルや対象に追加制限がある。
\dailyxsource{Anthropic (@AnthropicAI)}{https://x.com/AnthropicAI/status/2100646837799834096}

\section{Claude Projectsを並列クラウドスレッド型へ再設計}
Claude Projectsは、単なるフォルダから「司令塔となる1会話＋並列スレッド」へ再設計された。ベータは一部Pro/MaxのClaude Codeクラウドセッションから始まり、オフライン継続、共有メモリ、Overview、モバイルからの操作を掲げる。ローカルファイルや社内ネットワークへの対応は今後とされる。
\dailyxsource{Claude (@claudeai)}{https://x.com/claudeai/status/2100632677904744716}

\section{Claude、生体分子OSSモデル推論を最適化}
Anthropicは、Claudeが30超のオープンソース生体分子モデルの推論を最適化し、平均約4倍高速化したと発表した。GPU向けカスタムソフトを含むコードも公開し、Adaptyv Bioと5,000超のデザインを実験検証するコンペも実施する。速度値はAnthropic報告で、第三者再現は未確認である。
\dailyxsource{Anthropic (@AnthropicAI)}{https://x.com/AnthropicAI/status/2100701581109072332}

\section{Grok Botが音声対応}
xAIのGrok Bot公式は「Grok Bot can talk now」と発表し、デスクトップとモバイルへ数日かけて展開するとした。Elon Muskも音声対応を引用投稿した。基盤となる音声モデルや対応言語などの詳細は、観測された投稿本文では示されていない。
\dailyxsource{Grok Bot (@bot)}{https://x.com/bot/status/2100659463569170779}

\section{AlphaGenome Atlasの研究利用例を紹介}
Google DeepMindはAlphaGenome Atlasを用いた希少疾患、UK Biobank、調節配列研究の利用例を公式スレッドで紹介し、データベースを無料公開していると述べた。窓内投稿は利用事例の紹介であり、Atlas自体の初出発表かどうかは確定していない。
\dailyxsource{Google DeepMind (@GoogleDeepMind)}{https://x.com/GoogleDeepMind/status/2100586760036143330}

\section{World Labs Atlas、NVIDIA本社を3D探索}
World Labsは、Blackwell上で学習したAtlasへ32枚の画像を入力し、NVIDIA Voyager本社をリアルタイムに探索するデモを公開した。画像を3D空間文脈として新視点を生成する技術として紹介され、NVIDIA公式も引用した。一般ユーザーによる再現状況は未確認である。
\dailyxsource{World Labs (@theworldlabs)}{https://x.com/theworldlabs/status/2100631511703920763}

\section{Sam Altman、主要ローンチを来週へ延期}
Sam Altmanは、今週予定していた「一番楽しみにしていたローンチ」を来週へずらすと投稿した。引用元の「今週とDevDayで大きな出荷」という投稿は観測窓外であり、今回延期された具体的な製品名も特定されていない。Astra for Lawとの対応関係も不明である。
\dailyxsource{Sam Altman (@sama)}{https://x.com/sama/status/2100351958167220547}

\section{AltmanとHuangのホワイトハウス晩餐会出席を報道}
ABC Newsは関係者情報として、OpenAI CEO Sam AltmanとNVIDIA CEO Jensen Huangが、習近平氏を迎える来週のホワイトハウス晩餐会に出席予定だと報じた。観測窓内では本人やホワイトハウスによる公式出席確認はなく、報道ベースの情報として扱う。
\dailyxsource{ABC News (@ABC)}{https://x.com/ABC/status/2100640074513715275}

\section{NVIDIA、安全と能力の同時進展を強調}
NVIDIA公式は、Dumfries HouseでJensen Huangが国王や国際リーダーとAIの公益利用を議論したと報告し、「capability and safety must advance together」とするメッセージを紹介した。講演全文や共同声明は観測された投稿では確認されていない。
\dailyxsource{NVIDIA (@nvidia)}{https://x.com/nvidia/status/2100593782911582578}

\section{GPT-6 AstraのFrontierSWEスコア主張}
ProximalはFrontierSWEでGPT-6 Astraが65.5\%、Fable 5.1が56.3\%、GPT-5.6 Solが32.2\%だったと図付きで投稿した。ただしOpenAI公式による同ベンチマークの発表や公式リーダーボードは観測窓内で確認されておらず、未検証の主張として扱う。
\dailyxsource{Proximal (@ProximalHQ)}{https://x.com/ProximalHQ/status/2100705062448517317}

\section{a16z / fal、動画モデル高速化を議論}
a16zはfal創業者らとの対談を公開し、MiniMax H3の再実装でステップ削減とGPU利用率向上を行ったことや、Hollywoodが急成長セグメントであることを紹介した。35倍などの速度値は対談側の説明であり、独立した技術再現レポートは未確認である。
\dailyxsource{a16z (@a16z)}{https://x.com/a16z/status/2100655721905803571}
